% =============================================================================
% APPENDIX NC: FORMAL-NC: Non-Concavity, Fit, and the Landscape of Coherence
% =============================================================================
% Module: BCM2_07_FIT
% Category: FORMAL
% Version: 1.0 (January 2026)
% Status: Draft
% Language: EN
%
% This appendix formalizes the connection between complementarity and
% non-concavity, explaining why systems have multiple stable configurations
% (peaks) and why transitions between them require crossing valleys.
%
% =============================================================================

% =============================================================================
% CROSS-REFERENCE STRUCTURE
% =============================================================================
%
% CHAPTER LINKAGE (Required):
% - Primary Chapter: Chapter 7: Complementarity, Fit, and Non-Concavity
% - Secondary Chapters: Ch. 5 (Complementarity), Ch. 8 (Mathematical), Ch. 15 (HIERARCHY)
%
% This appendix provides foundation for:
%
% DEPENDENCIES (what this appendix REQUIRES):
% - Appendix B (CORE-HOW): Complementarity structure gamma
% - Appendix MIL: Milgrom-Roberts supermodularity theory
% - Appendix VII (FORMAL-GTC): General Theory of Complementarity
%
% DEPENDENTS (what REQUIRES this appendix):
% - Appendix IE (CORE-EIT): Emergent Intervention Theory (K vs Q in emergence)
% - Appendix VI (FORMAL-MEQ): Multiple Equilibria (landscape navigation)
% - Chapter 15: HIERARCHY emergence from complementarity
%
% =============================================================================

% -----------------------------------------------------------------------------
% APPENDIX SCOPE BOX
% -----------------------------------------------------------------------------

\begin{tcolorbox}[
    colback=orange!5!white,
    colframe=orange!75!black,
    title={\textbf{Appendix Scope: Ziel / In-Scope / Out-of-Scope / Constraints}},
    fonttitle=\bfseries
]

\textbf{Ziel (Objective):}
\begin{itemize}[nosep]
    \item Why does complementarity ($\gamma > 0$) create multiple stable configurations, and how do we distinguish internal coherence ($K$) from external fit ($Q$)?
\end{itemize}

\textbf{In-Scope:}
\begin{itemize}[nosep]
    \item Non-Concavity, Fit, and the Landscape of Coherence
\end{itemize}

\textbf{Out-of-Scope (delegiert an andere Appendices):}
\begin{itemize}[nosep]
    \item CORE-HOW $\rightarrow$ Appendix B
    \item FORMAL-GTC $\rightarrow$ Appendix VII
    \item CORE-EIT $\rightarrow$ Appendix IE
    \item FORMAL-MEQ $\rightarrow$ Appendix VI
    \item Milgrom-Roberts $\rightarrow$ Appendix MIL
\end{itemize}

\textbf{Constraints (Anwendungsgrenzen):}
\begin{itemize}[nosep]
    \item [TODO: Define when this appendix does NOT apply]
    \item [TODO: Define required prerequisites]
    \item [TODO: Define known limitations]
\end{itemize}

\textbf{Lieferobjekte (Deliverables):}
\begin{itemize}[nosep]
    \item 13 Equation(s)
    \item 4 Axiom(s)
    \item Worked Example(s)
\end{itemize}

\end{tcolorbox}


\section{Non-Concavity, Fit, and the Landscape of Coherence}
\label{app:formal-nc}

% -----------------------------------------------------------------------------
% HEADER BLOCK
% -----------------------------------------------------------------------------

\begin{tcolorbox}[colback=blue!5!white, colframe=blue!75!black,
    title=\textbf{Appendix: NC (FORMAL-NC)}]
\textbf{Category:} FORMAL (Formalization) \\
\textbf{Primary Chapter:} Chapter 7: Complementarity, Fit, and Non-Concavity \\
\textbf{Version:} 1.0 (January 2026) \\
\textbf{Core Contribution:} Formalizes why $\gamma > 0$ creates non-concave landscapes with multiple stable configurations, establishing the $K$ (coherence) vs.\ $Q$ (quality) distinction central to emergence theory.
\end{tcolorbox}

% -----------------------------------------------------------------------------
% CHAPTER LINKAGE BOX
% -----------------------------------------------------------------------------

\begin{tcolorbox}[colback=purple!5!white, colframe=purple!75!black,
    title=Chapter Linkage]

\textbf{Primary Chapter:} Chapter 7: Complementarity, Fit, and Non-Concavity
\begin{itemize}[nosep]
\item Section 7.1: Roberts' Insight $\leftrightarrow$ This Appendix Section NC.1
\item Section 7.2: Non-Concavity $\leftrightarrow$ This Appendix Section NC.2
\item Section 7.3: Fit as Coherence $\leftrightarrow$ This Appendix Section NC.3
\item Section 7.4: Valley Problem $\leftrightarrow$ This Appendix Section NC.4
\end{itemize}

\textbf{Secondary Chapters:}
\begin{itemize}[nosep]
\item Chapter 5: Complementarity --- provides $\gamma$ foundation
\item Chapter 8: Mathematical Foundations --- parameter calibration
\item Chapter 15: Decision Hierarchies --- HIERARCHY emerges from $\gamma$
\end{itemize}

\textbf{Reading Guide:}
\begin{itemize}[nosep]
\item For conceptual overview: Read Chapter 7 first
\item For formal supermodularity: See Appendix MIL (Milgrom-Roberts)
\item For emergence implications: See Appendix IE (CORE-EIT)
\end{itemize}

\end{tcolorbox}

% -----------------------------------------------------------------------------
% ABSTRACT
% -----------------------------------------------------------------------------

\begin{abstract}
This appendix formalizes how complementarity ($\gamma > 0$) creates non-concave
payoff landscapes with multiple stable configurations (peaks). We establish the
mathematical distinction between \textbf{coherence} ($K$)---internal fit among
elements---and \textbf{quality} ($Q$)---external fit with context. The ``valley
problem'' explains why organizational transformations fail: partial change moves
systems into valleys rather than toward new peaks. This framework connects
Chapter 7's organizational economics insights to the EBF emergence structure.
\end{abstract}

% -----------------------------------------------------------------------------
% QUICK REFERENCE BOX
% -----------------------------------------------------------------------------

\begin{tcolorbox}[colback=blue!5!white, colframe=blue!75!black,
    title=\textbf{Quick Reference: Non-Concavity and Fit}]

\textbf{Core Concepts:}
\begin{itemize}[nosep]
\item \textbf{Non-Concavity:} $\gamma > 0 \Rightarrow$ Hessian not negative semi-definite $\Rightarrow$ multiple local maxima
\item \textbf{Coherence ($K$):} Internal fit among configuration elements ($K \approx 1$ at peaks)
\item \textbf{Quality ($Q$):} External fit with context $\Psi$ (height of peak)
\item \textbf{Valley:} Region of low $\Pi$ between peaks (incoherent configurations)
\end{itemize}

\textbf{Key Equations:}
\begin{align}
\text{Supermodularity:} \quad & \frac{\partial^2 \Pi}{\partial x_i \partial x_j} > 0 \quad \forall i \neq j \label{eq:nc-supermod} \\
\text{Non-Concavity:} \quad & H = \nabla^2 \Pi \text{ has positive eigenvalue} \label{eq:nc-hessian} \\
\text{K-Q Relationship:} \quad & V(C, \Psi) = K(C) \cdot Q(C^*; \Psi) \label{eq:nc-kq}
\end{align}

\textbf{The Roberts Mapping:}
\begin{center}
\renewcommand{\arraystretch}{1.2}
\begin{tabular}{ll}
Roberts (2004) & EBF Framework \\
\hline
Fit & $K \approx 1$ \\
Misfit & $K < 1$ \\
Configuration quality & $Q(C^*; \Psi)$ \\
Multiple configurations & Multiple $C^*$ \\
\end{tabular}
\end{center}

\end{tcolorbox}

% =============================================================================
% FUNDAMENTAL QUESTION
% =============================================================================

\begin{tcolorbox}[colback=orange!5!white, colframe=orange!75!black,
    title=\textbf{Fundamental Question}]
\textbf{Why does complementarity ($\gamma > 0$) create multiple stable configurations, and how do we distinguish internal coherence ($K$) from external fit ($Q$)?}

This appendix answers: Given that elements of a system reinforce each other ($\gamma_{ij} > 0$), what are the mathematical consequences for the structure of possible configurations, and what does this imply for organizational change and behavioral interventions?
\end{tcolorbox}

% =============================================================================
% AXIOMS
% =============================================================================

\subsection{Axioms of Non-Concavity and Fit}
\label{sec:nc-axioms}

\begin{tcolorbox}[colback=gray!5!white, colframe=gray!75!black]
\textbf{Axiom NC-1 (Complementarity):}
For all pairs of configuration elements $i, j$:
\begin{equation}
\gamma_{ij} = \frac{\partial^2 \Pi}{\partial x_i \partial x_j} > 0 \quad \forall i \neq j
\label{ax:nc-1}
\end{equation}
Increasing one element raises the marginal return to others.

\textbf{Axiom NC-2 (Own-Concavity):}
Each element exhibits diminishing returns in isolation:
\begin{equation}
\frac{\partial^2 \Pi}{\partial x_i^2} < 0 \quad \forall i
\label{ax:nc-2}
\end{equation}

\textbf{Axiom NC-3 (Coherence Definition):}
Coherence $K$ measures proximity to the nearest local maximum:
\begin{equation}
K(C) = \frac{\Pi(C)}{\Pi(C^*_{local})} \in [0, 1]
\label{ax:nc-3}
\end{equation}
where $C^*_{local}$ is the nearest local maximum in configuration space.

\textbf{Axiom NC-4 (Quality Definition):}
Quality $Q$ measures the height of the current peak relative to the global maximum given context:
\begin{equation}
Q(C^*; \Psi) = \frac{\Pi(C^*; \Psi)}{\max_{C^*} \Pi(C^*; \Psi)}
\label{ax:nc-4}
\end{equation}
\end{tcolorbox}

% =============================================================================
% SECTION NC.1: THE MATHEMATICAL STRUCTURE OF NON-CONCAVITY
% =============================================================================

\subsection{NC.1: From Complementarity to Non-Concavity}
\label{sec:nc-nonconcavity}

\subsubsection{The Fundamental Theorem}

\begin{proposition}[Non-Concavity from Complementarity]
\label{prop:nc-main}
If the payoff function $\Pi(x_1, \ldots, x_n)$ exhibits complementarity:
\begin{equation}
\frac{\partial^2 \Pi}{\partial x_i \partial x_j} > 0 \quad \text{for all } i \neq j
\end{equation}
then $\Pi$ is \textbf{not globally concave}. The payoff landscape has multiple local maxima.
\end{proposition}

\begin{proof}
The Hessian matrix $H$ of $\Pi$ has structure:
\begin{equation}
H_{ij} = \begin{cases}
\frac{\partial^2 \Pi}{\partial x_i^2} < 0 & \text{if } i = j \text{ (own-concavity)} \\
\frac{\partial^2 \Pi}{\partial x_i \partial x_j} > 0 & \text{if } i \neq j \text{ (complementarity)}
\end{cases}
\end{equation}

For global concavity, $H$ must be negative semi-definite: all eigenvalues $\lambda_k \leq 0$.

Consider the vector $\mathbf{1} = (1, 1, \ldots, 1)^T$. Then:
\begin{equation}
\mathbf{1}^T H \mathbf{1} = \sum_i H_{ii} + \sum_{i \neq j} H_{ij}
\end{equation}

With sufficiently strong complementarity ($H_{ij} > |H_{ii}|/n$ for $i \neq j$), this sum is positive, implying at least one positive eigenvalue. Hence $H$ is indefinite, and $\Pi$ is not concave.
\end{proof}

\subsubsection{Interpretation: The Landscape Metaphor}

The non-concavity result means the ``landscape'' of $\Pi$ over configuration space has:

\begin{itemize}
\item \textbf{Peaks:} Local maxima---coherent configurations where all elements reinforce
\item \textbf{Valleys:} Local minima or saddles---incoherent configurations where elements clash
\item \textbf{Ridges:} Partial coherence---some elements fit, others don't
\end{itemize}

\begin{tcolorbox}[colback=yellow!5!white, colframe=yellow!75!black,
    title=\textbf{Intuition: Anna's Strategy Dilemma}]

Anna runs a consulting firm facing a strategic choice:

\textbf{Peak A (Efficiency):} Cost leadership, hierarchical structure, standardized processes, performance incentives

\textbf{Peak B (Innovation):} Differentiation, flat structure, flexible processes, team rewards

Both peaks are coherent ($K \approx 1$). But Anna cannot be ``half innovative''---that puts her in the valley:
\begin{itemize}
\item Innovation strategy + hierarchical structure = bureaucracy stifles ideas
\item Efficiency strategy + flat structure = chaos without coordination
\end{itemize}

The landscape has two peaks with a valley between. This is non-concavity in action.
\end{tcolorbox}

% =============================================================================
% SECTION NC.2: COHERENCE (K) VS. QUALITY (Q)
% =============================================================================

\subsection{NC.2: The K-Q Distinction}
\label{sec:nc-kq}

\subsubsection{Formal Definitions}

\begin{definition}[Coherence $K$]
\label{def:nc-coherence}
Coherence measures internal fit among configuration elements:
\begin{equation}
K(C) = \frac{\Pi(C)}{\Pi^{\max}_{local}} \in [0, 1]
\end{equation}
where $\Pi^{\max}_{local}$ is the nearest local maximum. $K \approx 1$ means the system is at a peak; $K < 1$ means it's in a valley or on a ridge.
\end{definition}

\begin{definition}[Quality $Q$]
\label{def:nc-quality}
Quality measures external fit with context:
\begin{equation}
Q(C^*; \Psi) = \frac{\Pi(C^*; \Psi)}{\Pi^{\max}_{global}(\Psi)}
\end{equation}
where $\Pi^{\max}_{global}(\Psi)$ is the height of the best peak given context $\Psi$.
\end{definition}

\subsubsection{The Critical Distinction}

\begin{tcolorbox}[colback=red!5!white, colframe=red!75!black,
    title=\textbf{Critical Insight: $K$ and $Q$ Are Independent}]

A system can have:
\begin{itemize}
\item \textbf{High $K$, High $Q$:} Coherent and well-adapted (ideal)
\item \textbf{High $K$, Low $Q$:} Coherent but obsolete (Kodak 2000)
\item \textbf{Low $K$, High $Q$:} Aiming for right peak but not there yet (in transition)
\item \textbf{Low $K$, Low $Q$:} Incoherent and maladapted (crisis)
\end{itemize}

\textbf{The Trap of Coherent Obsolescence:} High $K$ can mask low $Q$. A perfectly coherent firm can fail if context $\Psi$ shifts and its peak becomes a valley in the new landscape.
\end{tcolorbox}

\subsubsection{Mathematical Relationship}

The total value function decomposes:
\begin{equation}
V(C, \Psi) = K(C) \cdot Q(C^*; \Psi) \cdot \Pi^{\max}_{global}(\Psi)
\end{equation}

This shows:
\begin{itemize}
\item $K$ is about \textit{where you are relative to the nearest peak}
\item $Q$ is about \textit{which peak you're on relative to the best peak}
\item Both matter for total value $V$
\end{itemize}

% =============================================================================
% SECTION NC.3: THE VALLEY PROBLEM AND TRANSITIONS
% =============================================================================

\subsection{NC.3: The Valley Problem}
\label{sec:nc-valley}

\subsubsection{Why Partial Change Fails}

\begin{proposition}[Valley Crossing Requirement]
\label{prop:nc-valley}
To transition from Peak A to Peak B, a system must:
\begin{enumerate}
\item Leave Peak A (coherence falls: $K \downarrow$)
\item Cross the valley (period of low $\Pi$)
\item Climb to Peak B (coherence restored: $K \uparrow$)
\end{enumerate}
Partial change (moving some elements but not others) moves the system \textit{sideways into the valley} rather than toward the new peak.
\end{proposition}

\begin{equation}
\text{Transition path: } C^*_A \xrightarrow{K \downarrow} \text{Valley} \xrightarrow{K \uparrow} C^*_B
\end{equation}

\subsubsection{Why Sequential Change Is Worse}

\begin{center}
\renewcommand{\arraystretch}{1.2}
\begin{tabular}{lcccc}
\hline
\textbf{Change Strategy} & $\Delta x_1$ & $\Delta x_2$ & Direction & Result \\
\hline
Full transformation & $+0.8$ & $+0.8$ & Diagonal & Reach Peak B \\
Element 1 only & $+0.8$ & $0$ & Horizontal & Stuck in valley \\
Element 2 only & $0$ & $+0.8$ & Vertical & Stuck in valley \\
Sequential & $+0.8$, then $+0.8$ & & Two valleys & Worse than staying \\
\hline
\end{tabular}
\end{center}

The complementarity structure ($\gamma > 0$) means elements must change \textit{together}. Changing one without the others destroys coherence without gaining the benefits of the new configuration.

% =============================================================================
% SECTION NC.4: CONNECTION TO EIT AND EMERGENCE
% =============================================================================

\subsection{NC.4: Connection to Emergent Intervention Theory}
\label{sec:nc-eit}

\subsubsection{Why Non-Concavity Matters for Emergence}

The non-concavity framework connects directly to the 10C Emergence Classification (Appendix IE):

\begin{tcolorbox}[colback=green!5!white, colframe=green!75!black,
    title=\textbf{Connection to EIT: How $\gamma$ Creates Emergence}]

\textbf{From Chapter 7 to EIT:}
\begin{enumerate}
\item \textbf{HOW ($\gamma$) is Semi-Emergent:} Base complementarity $\gamma_0$ + context modulation $\gamma(\Psi)$
\item \textbf{Non-concavity from $\gamma > 0$:} Multiple stable configurations exist
\item \textbf{HIERARCHY emerges from $\gamma$:} Chapter 15 shows $N_{L2} = f(\gamma)$
\item \textbf{Peaks are emergent attractors:} $C^*$ emerges from the landscape structure
\end{enumerate}

\textbf{The Emergence Chain:}
\begin{equation}
\gamma > 0 \xrightarrow{\text{Prop. NC.1}} \text{Non-Concavity} \xrightarrow{\text{Multiple } C^*} \text{HIERARCHY Emergence}
\end{equation}

\textbf{K vs. Q in EIT Terms:}
\begin{itemize}
\item $K$ = coherence = how well the intervention elements fit together
\item $Q$ = quality = how well the configuration fits the behavioral context $\Psi$
\end{itemize}

An intervention portfolio can be internally coherent ($K \approx 1$) but poorly matched to the target segment ($Q < 1$)---or vice versa.
\end{tcolorbox}

\subsubsection{Implications for Intervention Design}

The non-concavity framework implies:

\begin{enumerate}
\item \textbf{Intervention portfolios must be coherent:} Combining T6 (Social) + T7 (Financial) creates complementarity conflict ($\gamma < 0$), pushing the portfolio into a valley.

\item \textbf{Partial interventions fail:} Addressing awareness (T1) without addressing willingness (T5) may leave the system worse off---in the valley between current behavior and target behavior.

\item \textbf{Context determines which peak is best:} The same intervention mix may be optimal in one $\Psi$ context and suboptimal in another.

\item \textbf{Transitions require sustained effort:} Moving from current behavior (Peak A) to target behavior (Peak B) requires crossing a valley---a period of discomfort and low payoff.
\end{enumerate}

% =============================================================================
% SECTION NC.5: WORKED EXAMPLES
% =============================================================================

\subsection{NC.5: Worked Examples}
\label{sec:nc-examples}

\subsubsection{Example 1: Kodak's Coherent Decline}

\textbf{Context:} Kodak in 2000 was highly coherent ($K \approx 1$):
\begin{itemize}
\item Strategy: Dominate film photography
\item Structure: Vertically integrated manufacturing
\item Processes: Optimized for chemical film production
\item Culture: Engineering excellence in film chemistry
\end{itemize}

\textbf{The Problem:} Digital photography shifted context $\Psi$. Kodak's peak (film excellence) became a valley in the new landscape.

\textbf{Analysis:}
\begin{align}
K_{\text{Kodak, 2000}} &\approx 1 \quad \text{(internally coherent)} \\
Q_{\text{Kodak, 2000}} &\approx 0.7 \quad \text{(film still profitable)} \\
Q_{\text{Kodak, 2010}} &\approx 0.1 \quad \text{(context shifted)}
\end{align}

High $K$ masked declining $Q$. Coherent obsolescence.

\subsubsection{Example 2: Digital Transformation Failure}

\textbf{Context:} A traditional manufacturer attempts digital transformation.

\textbf{The Wrong Way (Sequential):}
\begin{enumerate}
\item Year 1: Announce digital strategy ($\Delta s = +0.8$, $\Delta \sigma = 0$)
\item Result: In valley---strategy requires autonomy, structure prevents it
\item Years 2-5: Struggle in valley, key talent leaves
\item Year 6: Retreat to original configuration
\end{enumerate}

\textbf{The Right Way (Simultaneous):}
\begin{enumerate}
\item Year 0: Plan complete transformation (strategy, structure, culture, incentives)
\item Year 1: Change all elements together ($\Delta s = +0.8$, $\Delta \sigma = +0.8$, ...)
\item Years 2-3: Accept low $K$ during transition (cross valley)
\item Year 4: New coherent configuration emerges
\end{enumerate}

\textbf{Key Insight:} The simultaneous path is harder in the short run but reaches the new peak. The sequential path gets stuck in the valley.

% =============================================================================
% SECTION NC.6: CALIBRATION REQUIREMENTS
% =============================================================================
% RESULTS AND VALIDATION
% =============================================================================

\subsection{NC.6: Results and Empirical Validation}
\label{sec:nc-results}

The theoretical predictions of non-concavity have been validated empirically:

\begin{enumerate}
\item \textbf{Practice Bundling:} \citet{ichniowski1997} find that manufacturing plants adopt \textit{clusters} of complementary practices, not random combinations---consistent with multiple peaks.

\item \textbf{Management Practice Correlation:} \citet{bloom2007} document strong positive correlations among distinct management practices across 4,000+ firms in 12 countries---the signature of complementarity.

\item \textbf{Transformation Failure Rates:} ~70\% of organizational transformations fail \citep{mckinsey2015}, consistent with the valley problem: partial change lands in valleys.

\item \textbf{Coherent Decline:} Historical cases (Kodak, Blockbuster, Nokia) show high-$K$, declining-$Q$ patterns---coherent obsolescence as predicted.
\end{enumerate}

\textbf{Testable Predictions:}
\begin{itemize}[nosep]
\item Organizations with higher $\gamma$ (stronger complementarity) should show more clustered configurations
\item Successful transformations should exhibit simultaneous multi-element change
\item Context shifts ($\Delta\Psi$) should predict K-Q divergence before performance decline
\end{itemize}

% =============================================================================

\subsection{NC.7: Why Calibration Matters}
\label{sec:nc-calibration}

\subsubsection{What Cannot Be Derived from Principles}

The non-concavity framework provides \textbf{structure}:
\begin{itemize}
\item Landscapes have peaks and valleys
\item Complementarity creates multiple equilibria
\item Transitions require valley crossing
\end{itemize}

But the framework does \textbf{not} provide:
\begin{itemize}
\item How deep is the valley between Peak A and Peak B?
\item Which peak is highest given context $\Psi$?
\item How long can a system survive in the valley?
\end{itemize}

These are \textbf{quantitative questions} that require calibration from data---not derivation from principles. This connects to the WHERE dimension (Appendix BBB): parameters must be estimated, not assumed.

\begin{equation}
\boxed{\text{Framework: Structure} \quad + \quad \text{Calibration: Parameters} \quad = \quad \text{Prediction}}
\end{equation}

% =============================================================================
% SUMMARY BOX
% =============================================================================

\begin{tcolorbox}[colback=gray!10!white, colframe=gray!75!black,
    title=\textbf{Summary: Non-Concavity and Fit}]

\textbf{1. Complementarity Creates Non-Concavity:}
$\gamma > 0 \Rightarrow$ Hessian has positive eigenvalues $\Rightarrow$ multiple local maxima (peaks).

\textbf{2. K (Coherence) $\neq$ Q (Quality):}
\begin{itemize}[nosep]
\item $K$ = internal fit (am I at a peak?)
\item $Q$ = external fit (is my peak the right one for this context?)
\item High $K$, low $Q$ = coherent obsolescence (Kodak trap)
\end{itemize}

\textbf{3. The Valley Problem:}
Transitioning between peaks requires crossing a valley. Partial change gets stuck in the valley. Elements must change together.

\textbf{4. Connection to EIT:}
\begin{itemize}[nosep]
\item $\gamma > 0$ explains why HIERARCHY emerges (Chapter 15)
\item K vs. Q maps to intervention coherence vs. context fit
\item Non-concavity creates the landscape that emergence navigates
\end{itemize}

\textbf{5. Calibration Imperative:}
The framework provides structure (peaks, valleys, transitions). Parameters (valley depth, peak height) must be calibrated from data.

\end{tcolorbox}

% =============================================================================
% GLOSSARY OF SYMBOLS
% =============================================================================

\subsection{Glossary of Symbols}

For complete symbol definitions, see \textbf{Appendix GLS} (Central Glossary).

\begin{center}
\renewcommand{\arraystretch}{1.3}
\begin{tabular}{lll}
\hline
\textbf{Symbol} & \textbf{Definition} & \textbf{Reference} \\
\hline
$\gamma$ & Complementarity parameter & Appendix B (CORE-HOW) \\
$\Pi$ & Payoff function & Chapter 7 \\
$H$ & Hessian matrix of $\Pi$ & Eq. \ref{eq:nc-hessian} \\
$K$ & Coherence (internal fit) & Def. \ref{def:nc-coherence} \\
$Q$ & Quality (external fit) & Def. \ref{def:nc-quality} \\
$C^*$ & Equilibrium configuration & Appendix B \\
$\Psi$ & Context & Appendix V (CORE-WHEN) \\
\hline
\end{tabular}
\end{center}

% =============================================================================
% CRITICAL FOUNDATIONS
% =============================================================================

\subsection{Critical Foundations and Potential Objections}

\textbf{F1. Are real payoff landscapes truly non-concave?}

\textit{Response:} Empirical evidence from organizational economics (Ichniowski et al. 1997, Bloom \& Van Reenen 2007) shows that practices cluster---firms adopt bundles of complementary practices, not random combinations. This clustering is the signature of a non-concave landscape with multiple peaks.

\textbf{F2. Can the valley always be crossed?}

\textit{Response:} No. Some valleys are too deep (transition costs exceed resources), too wide (transition takes longer than the firm can survive), or have rising walls (context shifts during transition). The framework identifies these constraints; it does not guarantee successful transitions.

\textbf{F3. Is the K-Q distinction always meaningful?}

\textit{Response:} Yes, but the distinction is sharpest when context is changing. In stable contexts, high $K$ configurations tend to have high $Q$ (selection pressure). When $\Psi$ shifts, K-Q divergence becomes critical.

% =============================================================================
% OPEN ISSUES
% =============================================================================

\subsection{Open Issues and Research Agenda}

\begin{enumerate}
\item \textbf{Valley Depth Estimation:} Methods for calibrating transition costs between configurations.
\item \textbf{Dynamic Context:} How to model peak heights when $\Psi$ is changing during transition.
\item \textbf{Partial Coherence:} Can systems survive indefinitely on ridges, or do they inevitably slide into valleys or climb to peaks?
\item \textbf{Multi-Level Landscapes:} How do individual, organizational, and market landscapes interact?
\end{enumerate}

% =============================================================================
% REFERENCES
% =============================================================================

\section{References}
\label{sec:nc-references}

See \textbf{Master Bibliography}: \texttt{bcm2\_master\_references.bib}

\nocite{bcm_master}

Key references for this appendix:
\begin{itemize}[nosep]
\item \citet{roberts2004}: ``The Modern Firm'' --- organizational fit and complementarity
\item \citet{milgrom1990}: Supermodularity and strategic complementarity
\item \citet{ichniowski1997}: Empirical evidence on practice bundling
\item \citet{bloom2007}: Management practices and productivity
\end{itemize}

% =============================================================================
% CROSS-REFERENCE FOOTER
% =============================================================================

\begin{tcolorbox}[colback=gray!5!white, colframe=gray!50!black]
\textbf{Cross-References:}
\begin{itemize}[nosep]
\item \textbf{Depends on:} B (CORE-HOW), MIL (Milgrom-Roberts), VII (FORMAL-GTC)
\item \textbf{Required by:} IE (CORE-EIT), VI (FORMAL-MEQ), Chapter 15
\item \textbf{Related:} X (FORMAL-FOUND), XI (FORMAL-MODELS), BBB (CORE-WHERE)
\end{itemize}
\end{tcolorbox}
